\section{Kleines Evaluationsdesign}

\subsection{Ziel}

In kleiner Form soll geprueft werden, ob eine Textbeschreibung den Inhalt eines Graphviz-Workflows korrekt wiedergibt.

\subsection{Eingabe}

Die Eingabe besteht aus:

\begin{itemize}
  \item einem Workflow als \texttt{.gv}
  \item einer Textbeschreibung als \texttt{.txt}
\end{itemize}

\subsection{Was geprueft wird}

Folgende Eigenschaften werden geprueft:

\begin{itemize}
  \item Start wird sprachlich erwaehnt
  \item Ende wird sprachlich erwaehnt
  \item zentrale Aktivitaeten werden genannt
  \item OR-Splits werden als Alternative oder optionale Pfade erkennbar
  \item AND-Splits werden als parallel beschrieben
\end{itemize}

\subsection{Kleiner Score}

Der einfache Score besteht aus den folgenden Einzelmerkmalen:

\begin{itemize}
  \item \texttt{start\_mentioned}
  \item \texttt{end\_mentioned}
  \item \texttt{tasks\_coverage}
  \item \texttt{and\_parallel\_mentioned}
  \item \texttt{or\_alternative\_mentioned}
\end{itemize}

\subsection{Beispiel-Interpretation}

\begin{itemize}
  \item \texttt{5/5}: Beschreibung ist fuer einen ersten Prototyp sehr plausibel
  \item \texttt{3/5}: teilweise passend, aber wichtige Workflow-Eigenschaften fehlen
  \item \texttt{0--2/5}: Beschreibung ist als Validierungsobjekt unzureichend
\end{itemize}

\subsection{Limitationen}

\begin{itemize}
  \item reine Wortsuche ist noch keine echte semantische Validierung
  \item Synonyme koennen uebersehen werden
  \item falsche Reihenfolgen werden nur begrenzt erkannt
  \item spaeter sollte ein LLM- oder Graph-basierter Fehlabgleich ergaenzt werden
\end{itemize}
